\documentclass[11pt]{article}
\usepackage{amsmath,amssymb,booktabs,geometry}
\geometry{margin=1in}

\title{Step 199: Beurling--Nyman Harmonic Chain Extension to \(N=2000\)}
\author{six-birds-foundations-iii}
\date{}

\begin{document}
\maketitle

\section*{Target}

This step extends the Step 195 Beurling--Nyman harmonic-chain numerical
experiment.  For \(A_N=\{1/k:1\le k\le N\}\), the finite defect is
\[
  \delta_{A_N}^2=1-b_N^*G_N^\dagger b_N.
\]
The Gram entries use the arithmetic formula
\[
  G_{p,q}
  =
  \sum_{n\ge1}
  \frac{((n\bmod p)/p)((n\bmod q)/q)}{n(n+1)}.
\]

\section*{Method}

The scalar setup uses \texttt{mpmath} at 80 decimal digits.  The Gram matrix
is built by the Step 195 vectorized arithmetic series, truncated at
\[
  X_{\max}=40000,
\]
giving a per-entry tail bound \(2.5\times10^{-5}\).  The active matrix drops
the zero atom \(p=1\), then uses a symmetric eigen-pseudoinverse with
\[
  \mathrm{rcond}=10^{-12}.
\]

A full \(2000\times2000\) arbitrary-precision \texttt{mpmath} SVD was not
computationally feasible in this runtime.  Therefore this run is a
Step-195-comparable extension with an explicit \(N=1000\) calibration, not a
full high-precision dense-SVD certification.

\section*{Computed Data}

\begin{center}
\begin{tabular}{rrrrr}
\toprule
\(N\) & \(\delta_{A_N}^2\) & \(\delta_{A_N}^2\log N\) & rank & condition proxy\\
\midrule
1000 & 0.00646330050366617 & 0.0446468981738479 & 999 & \(5.0051\times10^6\)\\
1500 & 0.00611039863611640 & 0.0446866918788952 & 1499 & \(1.2654\times10^7\)\\
2000 & 0.00592901311818050 & 0.0450658503926355 & 1999 & \(2.4906\times10^7\)\\
\bottomrule
\end{tabular}
\end{center}

The \(N=1000\) calibration differs from the Step 195 value by
\[
  -8.09304602769102\times 10^{-5}
\]
in \(\delta^2\).  This is consistent with the larger truncation bound and the
changed active-matrix eigensolver.

\section*{Updated Decay Fit}

Combining inherited Step 195 values through \(N=1000\) with the new
\(N=1500,2000\) values gives:

\begin{center}
\begin{tabular}{llll}
\toprule
model & range & \(C\) & exponent / spread\\
\midrule
\(\delta^2\log N\) mean & \(N=200,500,1000,1500,2000\) & 0.0456632443937048 & spread 0.000926351\\
\(C/(\log N)^\alpha\) & \(N\ge20\) & 0.0547051524798759 & \(\alpha=1.096343229\)\\
recent \(C/(\log N)^\alpha\) & \(N\ge200\) & 0.0606633898663218 & \(\alpha=1.15047847\)\\
\bottomrule
\end{tabular}
\end{center}

\section*{Comparison with Step 195}

Step 195 reported a tail mean
\[
  \delta^2\log N\approx0.0467
\]
over \(N=80,\ldots,1000\), and a fit
\[
  \delta^2\sim C/(\log N)^\alpha,\qquad
  C\approx0.054,\quad \alpha\approx1.09.
\]
The extended data keep the same structure.  The updated all-tail fit gives
\[
  C\approx0.0547,\qquad \alpha\approx1.096.
\]
The direct values \(\delta^2\log N\) at \(N=1500,2000\) remain near
\[
  0.045.
\]

\section*{Verdict}

\[
\boxed{\texttt{V\_BN\_N2000\_consistent}.}
\]

The extension to \(N=2000\) confirms the slow-log decay pattern seen in Step
195, subject to the recorded truncation and dense-linear-algebra caveat.  It
is consistent with the Baez--Duarte RH-conditional Beurling--Nyman behavior.
It does not prove the infinite Beurling--Nyman closure or RH.

\end{document}

